\chapter{Modules over a Principal Ideal Domain}\label{ch:b3:modules}

Linear algebra over a \emph{ring} instead of a field: this small
change of hypothesis produces one of algebra's great unification
theorems. A module over $\Z$ is an abelian group; a module over
$K[X]$ is a vector space equipped with an endomorphism. The
structure theorem for finitely generated modules over a PID
therefore classifies, in one stroke, all finitely generated
abelian groups \emph{and} all endomorphisms up to similarity ---
Jordan's reduction, which Year 2 obtained by delicate inductions,
falls out as a corollary, together with its subtler sibling, the
\emph{rational} canonical form, valid over every field. The
computational engine is the Smith normal form, an arithmetic of
matrices worthy of Euclid.

Throughout, $A$ is a commutative ring, soon a PID; ``module''
means $A$-module.

\section{Modules, free modules}

\begin{definition}\label{def:b3:modules:module}
An \emph{$A$-module}\index{module} is an abelian group $(M, +)$
with a scalar multiplication $A \times M \to M$ satisfying the
vector-space axioms: $a(x + y) = ax + ay$, $(a + b)x = ax + bx$,
$(ab)x = a(bx)$, $1x = x$. Submodules, quotients $M/N$, morphisms
($A$-linear maps), direct sums $\bigoplus_i M_i$, and the
isomorphism theorems are defined and proved word for word as for
vector spaces and abelian groups; in particular $M/\ker f \cong
\operatorname{im} f$ for a morphism $f$.
\end{definition}

\begin{example}\label{ex:b3:modules:examples}
The three motivating cases.
\begin{enumerate}
  \item $A = K$ a field: modules are vector spaces.
  \item $A = \Z$: modules are exactly abelian groups ($nx$ is
        forced to be $x + \dots + x$), submodules are subgroups.
  \item $A = K[X]$: a module is a $K$-vector space $V$ together
        with the $K$-linear map $u\colon x \mapsto X\cdot x$ ---
        conversely, every pair $(V, u)$ with $u \in \mathcal L(V)$
        becomes a $K[X]$-module by $P \cdot x = P(u)(x)$. The
        submodules are precisely the $u$-stable subspaces.
\end{enumerate}
An ideal of $A$ is exactly a submodule of $A$; a quotient ring
$A/I$ is an $A$-module. Unlike vector spaces, modules can have
\emph{torsion}: in $\Z/6\Z$, the element $\bar 3 \ne 0$ is killed
by $2 \neq 0$.
\end{example}

\begin{definition}\label{def:b3:modules:free}
$M$ is \emph{finitely generated}\index{finitely generated module}
if $M = Ax_1 + \dots + Ax_n$ for some $x_i$. $M$ is
\emph{free}\index{free module} of rank $n$ if $M \cong A^n$,
i.e.\ if it has a \emph{basis} (a generating family that is
$A$-linearly independent). Every finitely generated $M$ is a
quotient of a free module: $(a_1, \dots, a_n) \mapsto \sum a_ix_i$
maps $A^n$ onto $M$.
\end{definition}

\begin{proposition}[Invariance of rank]\label{prop:b3:modules:rank}
If $A \neq 0$ and $A^m \cong A^n$, then $m = n$.
\end{proposition}

\begin{proof}
Pick a maximal ideal $\mathfrak m$ of $A$ (\cref{thm:b3:rings:krull})
and set $k = A/\mathfrak m$, a field. An isomorphism $f \colon A^m
\to A^n$ maps $\mathfrak m A^m$ into $\mathfrak m A^n$
(linearity), hence induces an isomorphism of quotients
\[
A^m/\mathfrak m A^m \;\cong\; A^n/\mathfrak m A^n,
\qquad\text{i.e.}\qquad k^m \cong k^n
\]
as $k$-vector spaces (the quotient $A^m/\mathfrak m A^m$ is
killed by $\mathfrak m$, so the $A$-action factors through $k$;
the images of the standard basis form a $k$-basis). Dimension
theory over the field $k$ gives $m = n$.
\end{proof}

\begin{theorem}[Submodules of free modules]\label{thm:b3:modules:submodule}
Let $A$ be a PID and $M \subseteq A^n$ a submodule. Then $M$ is
free of rank $\leq n$.
\end{theorem}

\begin{proof}
Induction on $n$. For $n = 1$: $M$ is an ideal, so $M = (0)$
(free of rank $0$) or $M = dA \cong A$ ($x \mapsto dx$ is
injective: domain). For $n > 1$: let $\pi \colon A^n \to A$ be the
last coordinate. Then $\pi(M)$ is an ideal, $(0)$ or $dA$. If
$(0)$: $M \subseteq A^{n-1} \times \{0\}$ and induction applies.
Otherwise choose $x_0 \in M$ with $\pi(x_0) = d$. Every $x \in M$
writes uniquely
\[
x = \underbrace{\frac{\pi(x)}{d}}_{\in A}\, x_0 +
\Bigl(x - \tfrac{\pi(x)}d x_0\Bigr),
\qquad x - \tfrac{\pi(x)}d x_0 \in M \cap \ker\pi
\]
($\pi(x) \in dA$, so the coefficient is in $A$). Thus $M = Ax_0
\oplus (M \cap \ker \pi)$: the sum is direct since $\pi(ax_0) =
ad = 0$ forces $a = 0$. By induction $M \cap \ker\pi \subseteq
\ker \pi \cong A^{n-1}$ is free of rank $\leq n - 1$; adjoining
$x_0$ (independent from $\ker\pi$ as just seen) gives a basis of
$M$ of cardinality $\leq n$.
\end{proof}

\begin{remark}\label{rem:b3:modules:presentation}
Consequently, over a PID every finitely generated module $M$ has a
\emph{finite presentation}: a surjection $\varphi\colon A^n \to M$
has free kernel with basis $c_1, \dots, c_k$ ($k \leq n$), and $M
\cong A^n / \operatorname{im}(C)$ where $C \in M_{n,k}(A)$ is the
matrix whose columns are the $c_j$. Understanding $M$ means
understanding a \emph{matrix over $A$} up to change of bases in
source and target --- the subject of the next section.
\end{remark}

\section{Smith normal form}

\begin{definition}\label{def:b3:modules:equivalence}
Two matrices $B, C \in M_{n,k}(A)$ are \emph{equivalent} if $C =
QBP$ with $Q \in GL_n(A)$, $P \in GL_k(A)$ (invertible
\emph{over $A$}: determinant in $A^\times$). Equivalent
presentation matrices define isomorphic modules $A^n/
\operatorname{im}$ (change bases in $A^n$ and $A^k$).
\end{definition}

\begin{theorem}[Smith normal form]\label{thm:b3:modules:smith}
Let $A$ be a PID and $B \in M_{n,k}(A)$. Then $B$ is equivalent to
a diagonal matrix\index{Smith normal form}
\[
\operatorname{diag}(d_1, d_2, \dots, d_r, 0, \dots, 0),
\qquad d_1 \mid d_2 \mid \cdots \mid d_r \neq 0 ,
\]
and the $d_i$ are unique up to associates: $d_1 \cdots d_i$ is a
gcd of the $i \times i$ minors of $B$ (in particular that gcd is
an invariant of equivalence). The $d_i$ are the \emph{invariant
factors}\index{invariant factors} of $B$.
\end{theorem}

\begin{proof}
\emph{Existence.} If $B = 0$, done. Otherwise, consider the set of
ideals $(b)$ generated by entries of matrices equivalent to $B$;
since $A$ is Noetherian, choose a matrix $B'$ equivalent to $B$
and an entry $d$ of $B'$ with $(d)$ \emph{maximal} in this set.
Move $d$ to position $(1,1)$ by row and column swaps.

\emph{Claim: $d$ divides every entry of $B'$.} First, column 1: if
$b_{i1}$ is not a multiple of $d$, let $e = \gcd(d, b_{i1}) = ud +
vb_{i1}$ (B\'ezout), so $(e) \supsetneq (d)$. The $2 \times 2$
matrix trick: acting on rows $1$ and $i$ by
\[
\begin{pmatrix} u & v \\ -\dfrac{b_{i1}}e & \dfrac de
\end{pmatrix}
\in GL_2(A)
\qquad
\Bigl(\det = \tfrac{ud + v b_{i1}}e = 1\Bigr)
\]
produces an equivalent matrix with entry $e$ in position $(1,1)$:
this contradicts maximality of $(d)$. So $d$ divides column $1$,
and symmetrically row $1$. Subtracting multiples of row 1 and
column 1 clears them: $B'$ is equivalent to $\begin{pmatrix} d &
0\\ 0 & B''\end{pmatrix}$. Next, $d$ divides every entry $b$ of
$B''$: add the row of $b$ to row 1 (an elementary operation; the
new first row contains $d$ and $b$-entries), and repeat the
column-clearing argument: a non-multiple would again improve
$(d)$. Now induct on the size: $B''$, all of whose entries are
divisible by $d$, has a Smith form $\operatorname{diag}(d_2,
\dots)$ whose entries remain divisible by $d$ (every entry of any
$QB''P$ is an $A$-combination of entries of $B''$); set $d_1 =
d$.

\emph{Uniqueness.} Let $D_i(B)$ denote a gcd of all $i \times i$
minors. Row and column operations, and more generally
multiplication by \emph{any} matrix, cannot shrink the gcd: the
$i \times i$ minors of $QB$ are $A$-combinations of those of $B$
(Cauchy--Binet expansion; or directly: each row of $QB$ is a
combination of rows of $B$, and minors are multilinear in rows).
So $D_i(QBP)$ and $D_i(B)$ divide each other: $D_i$ is an
equivalence invariant. On the diagonal form, the nonzero $i
\times i$ minors are the products of $i$ of the $d_j$'s, and
divisibility $d_1 \mid \dots \mid d_r$ makes $d_1 \cdots d_i$ the
gcd. Hence $d_1 \cdots d_i = D_i(B)$ up to units, and $d_i =
D_i/D_{i-1}$ is determined.
\end{proof}

\begin{method}\label{met:b3:modules:smith}
Over a Euclidean domain ($\Z$, $K[X]$), Smith reduction is an
algorithm --- no maximality argument needed: bring the entry of
smallest Euclidean size to position $(1,1)$; if it fails to divide
some entry of its row or column, a Euclidean division leaves a
strictly smaller remainder there --- swap it in and restart
(termination: sizes decrease); when it divides its whole row and
column, clear them; if it fails to divide an inner entry, add that
row to row $1$ and restart; recurse on the inner block. In
practice on integer matrices: compute $D_1 = \gcd$ of entries,
$D_2$, \dots\ via minors for small sizes, or run the algorithm.
\end{method}

\begin{example}[A Smith reduction, in full]
\label{ex:b3:modules:smithworked}
Reduce $M = \begin{pmatrix} 1 & 2 & 3\\ 4 & 5 & 6\\ 7 & 8 & 9
\end{pmatrix}$ over $\Z$. The corner $1$ divides everything:
clear its row and column ($L_2 \leftarrow L_2 - 4L_1$, $L_3
\leftarrow L_3 - 7L_1$, then $C_2 \leftarrow C_2 - 2C_1$, $C_3
\leftarrow C_3 - 3C_1$):
\[
M \sim \begin{pmatrix} 1 & 0 & 0\\ 0 & -3 & -6\\ 0 & -6 & -12
\end{pmatrix} .
\]
In the inner block, the corner $-3$ divides all entries: $L_3
\leftarrow L_3 - 2L_2$ and $C_3 \leftarrow C_3 - 2C_2$ clear it
to $\operatorname{diag}(-3, 0)$. Adjusting signs (multiply a
row by $-1$, a legal operation):
\[
M \sim \operatorname{diag}(1, 3, 0),
\qquad
\Z^3/M\Z^3 \cong \Z/3\Z \times \Z .
\]
Cross-check by determinantal divisors: $D_1 =
\gcd(\text{entries}) = 1$; every $2\times2$ minor of $M$ is a
multiple of $3$ (e.g.\ $\det\bigl(\begin{smallmatrix}1 & 2\\ 4
& 5\end{smallmatrix}\bigr) = -3$) and one equals $-3$: $D_2 =
3$; $D_3 = \det M = 0$. Hence $d_1 = 1$, $d_2 = 3$, $d_3 = 0$:
same answer. Two lessons: a zero invariant factor records the
rank drop (the cokernel picks up a free $\Z$ summand), and the
divisibility chain $1 \mid 3 \mid 0$ is the Smith certificate
--- a diagonal reduction that violates the chain (say
$\operatorname{diag}(2, 3)$, which the careless can produce
from $\bigl(\begin{smallmatrix}2 & 0\\ 0 &
3\end{smallmatrix}\bigr)$ by stopping too early: correct Smith
form $\operatorname{diag}(1, 6)$, as $D_1 = 1$ here!) is not
finished.
\end{example}

\begin{omfigure}
\begin{tikzpicture}[scale=0.72]
  \clip (-3.4,-1.3) rectangle (4.4,4.3);
  \foreach \x in {-3,...,4}
    \foreach \y in {-1,...,4}
      \fill[black!40] (\x,\y) circle (1.6pt);
  \fill[omDef!12, opacity=0.7] (0,0) -- (1,3) -- (3,3) -- (2,0)
    -- cycle;
  \foreach \a/\b in {0/0, 2/0, 4/0, -2/0, 1/3, 3/3, -1/3, -3/3,
                     2/6, -2/-6, 0/6}
    \fill[omThm] (\a,\b) circle (2.6pt);
  \draw[omThm, thick, ->] (0,0) -- (2,0)
    node[midway, below, font=\small] {$(2,0)$};
  \draw[omThm, thick, ->] (0,0) -- (1,3)
    node[midway, left, font=\small] {$(1,3)$};
\end{tikzpicture}

{\small The sublattice $L = \Z(2,0) + \Z(1,3)$ of $\Z^2$ (red
points). The Smith normal form of $\bigl(\begin{smallmatrix} 2 &
1\\ 0 & 3\end{smallmatrix}\bigr)$ is
$\operatorname{diag}(1, 6)$: in the adapted basis $f_1 = (1,3)$,
$f_2 = (0,1)$ of $\Z^2$, one has $L = \Z f_1 \oplus \Z\,6f_2$, so
$\Z^2/L \cong \Z/6\Z$ --- the index equals $\abs{\det} = 6$, the
area of the shaded fundamental domain.}
\end{omfigure}

\section{The structure theorem}

\begin{definition}\label{def:b3:modules:torsion}
Let $A$ be a domain and $M$ an $A$-module. The
\emph{torsion}\index{torsion} submodule is
\[
T(M) = \{x \in M : ax = 0 \text{ for some } a \neq 0\}
\]
(a submodule: if $ax = by = 0$ then $ab(x + y) = 0$, $ab \ne 0$).
$M$ is \emph{torsion-free} if $T(M) = 0$, a \emph{torsion module}
if $T(M) = M$.
\end{definition}

\begin{theorem}[Structure of finitely generated modules over a
PID]\label{thm:b3:modules:structure}
Let $A$ be a PID and $M$ a finitely generated $A$-module. There
exist a unique $r \in \N$ and nonzero nonunits $d_1 \mid d_2 \mid
\cdots \mid d_s$, unique up to associates, with\index{structure
theorem for modules}
\[
M \;\cong\; A^r \,\oplus\, A/(d_1) \oplus \cdots \oplus A/(d_s).
\]
Moreover $T(M) \cong A/(d_1)\oplus\dots\oplus A/(d_s)$ and $M/T(M)
\cong A^r$: a finitely generated \emph{torsion-free} module over a
PID is free.
\end{theorem}

\begin{proof}
\emph{Existence.} Present $M \cong A^n/\operatorname{im}(C)$
(\cref{rem:b3:modules:presentation}) and put $C$ in Smith form:
after the two base changes, $M \cong A^n / (d_1A \times \dots
\times d_rA \times 0 \times \dots \times 0) \cong A/(d_1) \oplus
\dots \oplus A/(d_r) \oplus A^{\,n - r}$. Discard the factors
where $d_i$ is a unit ($A/(d_i) = 0$); the divisibility chain
survives.

\emph{The torsion identification.} In the decomposition, $A^r$ is
torsion-free (a domain has no zero divisors) and each $A/(d_i)$
is torsion (killed by $d_i \ne 0$); a direct sum splits torsion
accordingly: $T(M) = \bigoplus_i A/(d_i)$ and $M/T(M) \cong A^r$.

\emph{Uniqueness of $r$}: $M/T(M) \cong A^r$ depends only on $M$,
and \cref{prop:b3:modules:rank} pins $r$ down.

\emph{Uniqueness of the $d_i$}: it suffices to treat the torsion
module $T = T(M)$. Decompose each $d_i$ into primes and split by
the Chinese remainder theorem (\cref{thm:b3:rings:crt}; distinct
primes generate comaximal ideals):
\[
A/(d) \cong \bigoplus_{p \mid d} A/\bigl(p^{v_p(d)}\bigr):
\qquad
T \cong \bigoplus_{p} \bigoplus_{j} A/\bigl(p^{k_{p,j}}\bigr),
\]
the \emph{elementary divisors}\index{elementary divisors}
$p^{k_{p,j}}$. Conversely the $d_i$ are reconstructed from the
multiset of elementary divisors ($d_s$ = product of the highest
power of each prime, etc.), so it suffices to prove the
multiset $\{k_{p,j}\}_j$ is determined by $T$, for each prime
$p$. Fix $p$; for $j \geq 1$ consider the $A/(p)$-vector spaces
$p^{j-1}T/p^jT$. On a cyclic factor $A/(p^k)$:
\[
p^{j-1}\bigl(A/(p^k)\bigr)\big/p^{j}\bigl(A/(p^k)\bigr)
\cong
\begin{cases}
A/(p) & \text{if } j \leq k,\\
0 & \text{if } j > k,
\end{cases}
\]
and on a factor $A/(q^k)$, $q \neq p$: multiplication by $p$ is
bijective there ($p$ invertible mod $q^k$: B\'ezout), so the
quotient is $0$. Direct sums pass through: $\dim_{A/(p)}
p^{j-1}T/p^jT = \#\{i : k_{p,i} \geq j\}$. These intrinsic
dimensions determine the multiset of exponents.
\end{proof}

\begin{corollary}[Finitely generated abelian groups]
\label{cor:b3:modules:abelian}
Every finitely generated abelian group is $\Z^r \times
\Z/d_1\Z\times\dots\times\Z/d_s\Z$ with $d_1 \mid \dots \mid
d_s$, uniquely. Every \emph{finite} abelian group is a product of
cyclic groups of prime-power order, unique as a
multiset.\index{abelian group, structure}
\end{corollary}

\begin{example}\label{ex:b3:modules:abelian}
The abelian groups of order $p^n$ correspond to \emph{partitions}
of $n$: for $p^4$: $\Z/p^4$, $\Z/p^3\times\Z/p$,
$(\Z/p^2)^2$, $\Z/p^2 \times (\Z/p)^2$, $(\Z/p)^4$ --- five
groups, as $4$ has five partitions. Mixed orders multiply the
counts prime by prime (CRT): there are $5 \times 2$ abelian
groups of order $2^4 \cdot 3^2 = 144$.
\end{example}

\section{Application: canonical forms of endomorphisms}

Let $K$ be a field, $V$ a $K$-vector space of finite dimension
$n$, and $u \in \mathcal L(V)$; make $V$ a $K[X]$-module via $P
\cdot x = P(u)(x)$ (\cref{ex:b3:modules:examples}). This module
is finitely generated (a $K$-basis generates) and torsion: for
each $x$, the $n+1$ vectors $x, u(x), \dots, u^n(x)$ are
$K$-dependent, providing a nonzero annihilating polynomial.

\begin{definition}\label{def:b3:modules:companion}
For $P = X^m + a_{m-1}X^{m-1} + \dots + a_0$ monic, the
\emph{companion matrix}\index{companion matrix} is
\[
C_P = \begin{pmatrix}
0 & & & -a_0\\
1 & \ddots & & -a_1\\
& \ddots & 0 & \vdots\\
& & 1 & -a_{m-1}
\end{pmatrix} :
\]
the matrix of ``multiplication by $X$'' on $K[X]/(P)$ in the
basis $1, \bar X, \dots, \bar X^{m-1}$.
\end{definition}

\begin{theorem}[Frobenius: rational canonical form]
\label{thm:b3:modules:frobenius}
There is a unique sequence of monic nonconstant polynomials $P_1
\mid P_2 \mid \dots \mid P_s$ (the \emph{similarity
invariants}\index{similarity invariants} of $u$) such that, as
$K[X]$-modules,
\[
V \cong K[X]/(P_1) \oplus \cdots \oplus K[X]/(P_s):
\]
in a suitable basis, $u$ has block-diagonal matrix
$\operatorname{diag}(C_{P_1}, \dots, C_{P_s})$. Moreover:
\begin{enumerate}
  \item $P_s = \mu_u$ (minimal polynomial) and $P_1\cdots P_s =
        \chi_u$ (characteristic polynomial); in particular $\mu_u
        \mid \chi_u$ \emph{(Cayley--Hamilton re-proved)} and
        $\chi_u \mid \mu_u^{\,s}$, so $\chi_u$ and $\mu_u$ have
        the same irreducible factors.
  \item Two endomorphisms (or square matrices) are similar iff
        they have the same similarity invariants.
\end{enumerate}
\end{theorem}

\begin{proof}
The structure theorem (\cref{thm:b3:modules:structure}) applied
to the PID $K[X]$: the torsion module $V$ decomposes with
invariant factors $P_i$, normalized monic (units of $K[X]$ are
$K^\times$); no free part occurs ($V$ is torsion). On each cyclic
factor $K[X]/(P_i)$, multiplication by $X$ has matrix $C_{P_i}$
in the basis of powers of $\bar X$: concatenating bases gives the
block form.

(1) The annihilator of $V = \bigoplus K[X]/(P_i)$ is $(P_1)\cap
\dots\cap(P_s) = (P_s)$ (divisibility chain: $P_s$ is a common
multiple, and the class of $1$ in the last factor is killed
exactly by $(P_s)$): $\mu_u = P_s$. For $\chi_u$: on a cyclic
factor, $\chi_{C_P} = P$, by induction on $m = \deg P$. Expanding
$\det(XI_m - C_P)$ along the first row (whose entries are $X$, then
zeros, then $a_0$ in the last column):
\[
\det(XI_m - C_P) = X\,\det\bigl(XI_{m-1} - C_{\tilde P}\bigr)
+ (-1)^{1+m}\,a_0\,\det L ,
\]
where $\tilde P = X^{m-1} + a_{m-1}X^{m-2} + \dots + a_1$ (same
shape, one size down) and $L$ is triangular with diagonal
$(-1, \dots, -1)$, so $\det L = (-1)^{m-1}$. By induction the
first term is $X\tilde P$, and the second is $a_0$: the total is
$X\tilde P + a_0 = P$ (base case $m=1$: $\det(X + a_0) = P$).
Determinants multiply over blocks:
$\chi_u = \prod P_i$. Cayley--Hamilton: $\chi_u \in (\mu_u)$
since $P_s \mid$ each\dots\ conversely each $P_i \mid P_s$, so
$\chi_u = \prod P_i$ divides $P_s^{\,s} = \mu_u^s$; and $\mu_u =
P_s$ divides $\chi_u$ as one of its factors.

(2) Similar endomorphisms are conjugate module structures, hence
have equal invariants (uniqueness in
\cref{thm:b3:modules:structure}); conversely equal invariants
give isomorphic $K[X]$-modules, and a module isomorphism is
exactly a linear bijection intertwining the two endomorphisms: a
similarity.
\end{proof}

\begin{corollary}[Similarity is insensitive to field extension]
\label{cor:b3:modules:descent}
Let $K \subseteq L$ be fields and $M, N \in M_n(K)$. If $M$ and
$N$ are similar over $L$, they are similar over $K$.
\end{corollary}

\begin{proof}
The similarity invariants of $M$ are computed by Smith's minor
formula (\cref{thm:b3:modules:smith}) applied to the presentation
matrix $XI_n - M$ over $K[X]$ --- indeed the $K[X]$-module
$V_M = K^n$ has presentation $XI_n - M$: the map $K[X]^n \to V_M$,
$(Q_i) \mapsto \sum Q_i(M)e_i$, is onto with kernel generated by
the columns of $XI_n - M$ (a direct verification: modulo those
columns, every element of $K[X]^n$ reduces to a constant vector,
and constant vectors map bijectively; the weekend problem spells
this out). Gcds of polynomials do not change under field
extension: if $d$ is the monic gcd in $K[X]$ of a family $(f_j)$,
B\'ezout gives $d = \sum u_jf_j$ with $u_j \in K[X]$, so every
common divisor of the $f_j$ \emph{in $L[X]$} divides $d$; as $d$
is itself a common divisor, it is the gcd in $L[X]$ too. Hence
the invariant factors of $XI_n - M$, quotients of successive
minor gcds, are the same over $K$ and over $L$: $M, N$ have the
same similarity invariants over $L$ iff over $K$; conclude by
\cref{thm:b3:modules:frobenius}(2).
\end{proof}

\begin{theorem}[Jordan form, re-derived]\label{thm:b3:modules:jordan}
Suppose $\chi_u$ splits over $K$ (e.g.\ $K = \C$). Applying to
$V$ the elementary-divisor decomposition (proof of
\cref{thm:b3:modules:structure}) instead of invariant factors:
\[
V \cong \bigoplus_{\lambda, j} K[X]\big/\bigl((X -
\lambda)^{k_{\lambda,j}}\bigr),
\]
and in the basis $\bigl(\overline{(X-\lambda)^{k-1}}, \dots,
\overline{(X - \lambda)}, \bar 1\bigr)$ of each factor, $u$ acts
as the Jordan block $J_k(\lambda)$: every endomorphism with split
characteristic polynomial has a Jordan basis, and the multiset of
blocks $(\lambda, k)$ is unique.\index{Jordan form}
\end{theorem}

\begin{proof}
The elementary divisors of the torsion module $V$ are the
$(X - \lambda)^k$ with $X - \lambda$ ranging over the irreducible
factors of $\mu_u$ (which splits, since $\chi_u$ does and both
have the same irreducible factors,
\cref{thm:b3:modules:frobenius}). In $W = K[X]/((X-\lambda)^k)$,
put $f_j = \overline{(X - \lambda)^{k-j}}$ for $j = 1, \dots, k$:
then $(X - \lambda)f_j = f_{j-1}$ (with $f_0 = 0$), i.e.\ $u(f_j)
= \lambda f_j + f_{j-1}$: the matrix of $u$ on $(f_1, \dots,
f_k)$ is exactly $J_k(\lambda)$ (ones above the diagonal).
Uniqueness of the multiset of elementary divisors is
\cref{thm:b3:modules:structure}.
\end{proof}

\begin{remark}\label{rem:b3:modules:overview}
The hierarchy of canonical forms is now transparent: the
\emph{rational} form exists over every field and detects
similarity absolutely (\cref{cor:b3:modules:descent}); the
\emph{Jordan} form is its refinement when $\chi_u$ splits.
Year 2's dimension-counting proofs of Jordan's theorem are
subsumed: all the combinatorics was the arithmetic of the PID
$K[X]$.
\end{remark}

\section{Exercises}

\begin{exercise}[$\star$]\label{exo:b3:modules:1}
(a) Show that $\Q$ is not finitely generated as a $\Z$-module.
(b) Show that $\Q$ is torsion-free but not free.
(c) Why does neither statement contradict
\cref{thm:b3:modules:structure}?
\end{exercise}

\begin{exercise}[$\star$]\label{exo:b3:modules:2}
List the abelian groups of order $360$ up to isomorphism, in both
elementary-divisor and invariant-factor forms. How many abelian
groups of order $p^5$ are there?
\end{exercise}

\begin{exercise}[$\star$]\label{exo:b3:modules:3}
Compute the Smith normal form over $\Z$ of
\[
B = \begin{pmatrix} 2 & 4\\ 6 & 8 \end{pmatrix},
\qquad
C = \begin{pmatrix} 2 & 0 & 0\\ 0 & 3 & 0\\ 0 & 0 & 12
\end{pmatrix},
\]
and identify the abelian groups $\Z^2/B\Z^2$ and $\Z^3/C\Z^3$.
\end{exercise}

\begin{exercise}[$\star\star$]\label{exo:b3:modules:4}
Let $L \subseteq \Z^n$ be a subgroup of rank $n$ with basis the
columns of $B \in M_n(\Z)$, $\det B \neq 0$. Show that
$\Z^n/L$ is finite of cardinality $\abs{\det B}$, and that $\Z^n/L
\cong \prod_i \Z/d_i\Z$ for the invariant factors $d_i$ of $B$.
Illustrate with $L = \Z(2,0) + \Z(1,3)$.
\end{exercise}

\begin{exercise}[$\star\star$]\label{exo:b3:modules:5}
Let $A$ be a domain. (a) Verify that $T(M)$ is a submodule and
that $M/T(M)$ is torsion-free. (b) Show that the ideal $(X, Y)$
of $K[X,Y]$, as a $K[X,Y]$-module, is torsion-free but not free:
the structure theorem genuinely needs the PID hypothesis.
\end{exercise}

\begin{exercise}[$\star\star$]\label{exo:b3:modules:6}
(a) Show that $2\Z$ has no direct complement in the $\Z$-module
$\Z$: submodules of free modules are free
(\cref{thm:b3:modules:submodule}), but direct summands they need
not be. (b) Show that if $M \subseteq A^n$ ($A$ a PID) satisfies:
$A^n/M$ is torsion-free, then $M$ \emph{is} a direct summand.
\end{exercise}

\begin{exercise}[$\star\star$]\label{exo:b3:modules:7}
Solve in $\Z^2$ the system
\[
\begin{cases}
2x + 4y \equiv b_1 \pmod{20},\\
6x + 8y \equiv b_2 \pmod{20},
\end{cases}
\]
for which pairs $(b_1, b_2)$ solutions exist, using the Smith
form of \cref{exo:b3:modules:3} (invertible changes of variables
on both sides).
\end{exercise}

\begin{exercise}[$\star\star$]\label{exo:b3:modules:8}
(a) Determine all similarity invariants and possible Jordan forms
of a nilpotent $4 \times 4$ matrix, sorted by the partition of
$4$ they realize.
(b) Exhibit two $4\times4$ complex matrices with the same
characteristic \emph{and} minimal polynomials that are not
similar, and prove that for $n \leq 3$ this cannot happen.
\end{exercise}

\begin{exercise}[$\star\star\star$]\label{exo:b3:modules:9}
Let $u \in \mathcal L(V)$, $\dim V = n$. Show that the following
are equivalent: (i) $V$ is a cyclic $K[X]$-module (there is $x$
with $V = K[u]x$, a \emph{cyclic vector}); (ii) $\mu_u = \chi_u$;
(iii) $s = 1$ in \cref{thm:b3:modules:frobenius}. Deduce that a
companion matrix has a cyclic vector, and determine when a
diagonal matrix has one.
\end{exercise}

\begin{exercise}[$\star\star\star$]\label{exo:b3:modules:10}
For $M \in M_n(\Z)$ viewed as an endomorphism of $\Z^n$, prove
the index formula: if $\det M \ne 0$, then
$[\Z^n : M\Z^n] = \abs{\det M}$, and deduce that $M \in GL_n(\Z)$
iff $\det M = \pm 1$. Application: the group $\Z^2$ has exactly
$\sigma_1(m) = \sum_{d \mid m} d$ subgroups of index $m$.
\emph{(Count matrices in Hermite form $\bigl(\begin{smallmatrix}
a & b\\ 0 & d\end{smallmatrix}\bigr)$, $ad = m$, $0 \leq b <
d$.)}
\end{exercise}

\begin{exercise}[$\star\star$]\label{exo:b3:modules:11}
(Equations $x^k = e$ in abelian groups) Let $G$ be a finite
abelian group with invariant factors $d_1 \mid d_2 \mid \dots
\mid d_s$.
(a) Show that for every $k \geq 1$,
\[
\#\{x \in G : x^k = e\} \;=\; \prod_{i=1}^{s}\gcd(k, d_i) .
\]
(b) Deduce: a finite abelian group is cyclic if and only if
for every $k$, the equation $x^k = e$ has at most $k$
solutions.
(c) Recover the cyclicity of finite subgroups of $K^\times$
($K$ a field, \cref{ch:b3:galois}): why does the polynomial
$X^k - 1$ guarantee the criterion of (b)?
\end{exercise}

\begin{exercise}[$\star\star\star$]\label{exo:b3:modules:12}
(Elementary matrices generate) (a) Show that $M \in M_n(\Z)$
is invertible in $M_n(\Z)$ iff $\det M = \pm1$.
(b) Show that $SL_2(\Z)$ is generated by the two elementary
matrices
$E = \bigl(\begin{smallmatrix}1 & 1\\ 0 &
1\end{smallmatrix}\bigr)$ and
$F = \bigl(\begin{smallmatrix}1 & 0\\ 1 &
1\end{smallmatrix}\bigr)$.
\emph{(Run the Euclidean algorithm on the first column of $M
\in SL_2(\Z)$ by left multiplications by powers of $E, F$,
reaching $\pm\bigl(\begin{smallmatrix}1 & *\\ 0 &
1\end{smallmatrix}\bigr)$; finish by hand --- note
$-I = (EF^{-1}E)^2$.)}
(c) Explain the connection with Smith reduction: over $\Z$,
row and column operations of determinant $1$ suffice to
diagonalize, up to signs.
\end{exercise}

\section{Problem: the commutant and the double commutant}

\begin{problem}[{Weekend problem --- rational form, commutant,
bicommutant}]\label{pb:b3:modules:1}
Let $K$ be a field, $V$ a $K$-vector space of dimension $n \geq
1$, and $u \in \mathcal L(V)$. We study the \emph{commutant}
\[
\mathcal C(u) = \{v \in \mathcal L(V) : uv = vu\},
\]
a subalgebra of $\mathcal L(V)$ containing $K[u] = \{P(u) : P \in
K[X]\}$, and we prove Frobenius' dimension formula and the double
commutant theorem: $\mathcal C(\mathcal C(u)) = K[u]$. Throughout,
$V$ is the $K[X]$-module defined by $u$, with invariant factors
$P_1 \mid \cdots \mid P_s$ and cyclic decomposition $V =
\bigoplus_{i=1}^s V_i$, $V_i = K[u]\,x_i \cong K[X]/(P_i)$, $n_i
= \deg P_i$ (\cref{thm:b3:modules:frobenius}).

\medskip
\noindent\textbf{Part I --- The presentation matrix $XI - M$, and
warm-ups.}
\begin{enumerate}
  \item Let $M \in M_n(K)$ and let $\varphi \colon K[X]^n \to V_M
        = K^n$ send $(Q_1, \dots, Q_n)$ to $\sum_i Q_i(M)e_i$.
        Show that $\varphi$ is a surjective morphism of
        $K[X]$-modules and that every column of $XI_n - M$ lies
        in $\ker\varphi$.
  \item Show that, modulo the columns of $XI_n - M$, every
        element of $K[X]^n$ is congruent to a constant vector
        \emph{(reduce degrees using $Xe_i \equiv \sum_j m_{ji}e_j$)},
        and deduce $\ker\varphi = (XI_n - M)\,K[X]^n$: the module
        $V_M$ has presentation matrix $XI_n - M$. Recover
        \cref{cor:b3:modules:descent}'s starting point: the
        similarity invariants of $M$ are the nonunit invariant
        factors of $XI_n - M$.
  \item Compute the similarity invariants of: a scalar matrix
        $\lambda I_n$; a diagonal matrix with distinct diagonal
        entries; the $n \times n$ Jordan block $J_n(0)$;
        $\operatorname{diag}(J_2(0), J_1(0))$ for $n = 3$.
  \item Show that $\dim K[u] = \deg \mu_u = n_s$.
\end{enumerate}

\medskip
\noindent\textbf{Part II --- Morphisms between cyclic modules.}
\begin{enumerate}[resume]
  \item Let $P, Q$ be monic nonconstant. Show that a
        $K[X]$-morphism $f \colon K[X]/(P) \to K[X]/(Q)$ is
        determined by $f(\bar 1)$, and that $\bar c \in K[X]/(Q)$
        can serve as $f(\bar 1)$ iff $P\bar c = 0$ in $K[X]/(Q)$.
  \item Deduce
        \[
        \operatorname{Hom}_{K[X]}\bigl(K[X]/(P),\, K[X]/(Q)\bigr)
        \;\cong\; K[X]\big/\bigl(\gcd(P, Q)\bigr),
        \]
        of dimension $\deg \gcd(P, Q)$ over $K$. \emph{(Show that
        the solutions $\bar c$ of $P\bar c = 0$ in $K[X]/(Q)$
        form the cyclic submodule generated by
        $\overline{Q/\gcd(P,Q)}$.)}
  \item Prove \emph{Frobenius' formula}:
        \[
        \dim_K \mathcal C(u) = \sum_{i,j=1}^{s} \deg\gcd(P_i,
        P_j) = \sum_{i=1}^{s} (2s - 2i + 1)\, n_i .
        \]
        \emph{(A commuting $v$ is exactly a $K[X]$-endomorphism of
        $V$; decompose $\operatorname{End}(\bigoplus_i V_i)$ as
        matrices of morphisms $V_j \to V_i$ and use the
        divisibility chain.)}
  \item Deduce $\dim \mathcal C(u) \geq n$, with equality iff $u$
        is cyclic ($s = 1$), and compute $\dim\mathcal C(u)$ for
        $u = \lambda\,\mathrm{id}$: both extremes of the formula.
  \item Verify Frobenius' formula directly for
        $\operatorname{diag}(J_2(0), J_1(0))$ by computing the
        commutant explicitly as $3\times3$ matrices.
\end{enumerate}

\medskip
\noindent\textbf{Part III --- The double commutant theorem.}
Let $w \in \mathcal C(\mathcal C(u))$; we prove $w \in K[u]$.
\begin{enumerate}[resume]
  \item Show $K[u] \subseteq \mathcal C(\mathcal C(u))$, and that
        every $w \in \mathcal C(\mathcal C(u))$ commutes with $u$
        --- so the inclusion to be proved,
        $\mathcal C(\mathcal C(u)) \subseteq K[u]$, is a genuine
        sharpening of $w \in \mathcal C(u)$.
  \item Suppose first that $u$ is \emph{cyclic}, $V = K[u]x$.
        Show directly that $\mathcal C(u) = K[u]$
        \emph{(evaluate a commuting $v$ on $x$: $v(x) = P(u)x$
        for some $P$, and compare $v$ with $P(u)$ on the basis
        $u^k x$)}, and conclude the theorem in this case.
  \item Back to the general case. For each $i$, let $\pi_i\colon
        V \to V_i$ be the projection along the other summands.
        Show $\pi_i \in \mathcal C(u)$, and deduce that $w$
        preserves each $V_i$ and commutes with $u_i =
        u\restriction_{V_i}$; conclude via question 11 applied to
        the cyclic $u_i$: there are polynomials $Q_i$ with
        $w\restriction_{V_i} = Q_i(u)\restriction_{V_i}$.
  \item It remains to glue the $Q_i$ into one polynomial. For $i
        \leq j$ (so $P_i \mid P_j$), show that $\eta_{ij} \colon
        V_j \to V_i$, $R(u)x_j \mapsto R(u)x_i$, is a
        well-defined $K[X]$-morphism \emph{(what must be checked
        is that $R(u)x_j = 0$ implies $R(u)x_i = 0$)}, and that
        $\tilde\eta_{ij} = \eta_{ij}\circ\pi_j$, extended by $0$
        on the other summands, lies in $\mathcal C(u)$.
  \item Using $w\tilde\eta_{ij} = \tilde\eta_{ij}w$, show $Q_i
        \equiv Q_j \pmod{P_i}$ for $i \leq j$. Deduce that $Q =
        Q_s$
        satisfies $Q \equiv Q_i \pmod {P_i}$ for all $i$, hence
        $w = Q(u)$ on every $V_i$, hence on $V$:
        \[
        \boxed{\ \mathcal C(\mathcal C(u)) = K[u].\ }
        \]
  \item (Coda) Deduce from the theorem: if $v$ commutes with
        \emph{every} matrix commuting with $u$, and $u$ is
        cyclic, then $v$ is a polynomial in $u$; and give an
        example showing $\mathcal C(u) = K[u]$ fails for $u =
        \mathrm{id}$, $n \geq 2$ --- where exactly does
        cyclicity enter?
\end{enumerate}

\medskip
\noindent\textbf{Part IV --- Dividends of the similarity
invariants.} The rational canonical form is a machine; here are
five of its classical outputs.
\begin{enumerate}[resume]
  \item (Transpose) Show that every $M \in M_n(K)$ is similar to
        its transpose ${}^tM$. \emph{(The operations that bring
        $XI - M$ to Smith form, transposed, bring $XI - {}^tM$ to
        the same Smith form: equal similarity invariants.)}
  \item (Descent of similarity) Let $K \subseteq L$ be a field
        extension and $M, N \in M_n(K)$. Show that if $M$ and $N$
        are similar over $L$, they are similar over $K$.
        \emph{(The Smith form of $XI - M$ computed in $K[X]$ is
        still a Smith form in $L[X]$ --- why do the invariant
        factors not change?)} Consequence worth memorizing: two
        real matrices conjugate in $GL_n(\C)$ are conjugate in
        $GL_n(\R)$.
  \item (Nilpotent classification) Let $u$ be nilpotent. Show
        that the number of blocks of size $\geq k$ in its
        decomposition into nilpotent Jordan blocks equals
        $\operatorname{rk}u^{k-1} - \operatorname{rk}u^k$, and
        deduce: nilpotent classes of $M_n(K)$, for any field $K$,
        are in bijection with the partitions of $n$. How many
        nilpotent classes in $M_5(K)$?
  \item (A concrete pair) Determine the similarity invariants of
        the derivation $D\colon P \mapsto P'$ acting on the space
        $K_{n-1}[X]$ of polynomials of degree $< n$: (a) for $K =
        \Q$; (b) for $K = \mathbb F_p$ with $p < n$
        \emph{(in characteristic $p$, $(X^p)' = 0$: compute
        $\ker D^k$ and use question 18)}.
  \item (Conjugacy classes of $GL_2(\mathbb F_q)$) Using
        invariant factors, show that every class of
        $GL_2(\mathbb F_q)$ is of exactly one of four types:
        central $aI$; diagonalizable with two distinct
        eigenvalues $a \neq b$ in $\mathbb F_q^\times$;
        non-semisimple with minimal polynomial $(X - a)^2$;
        cyclic with irreducible characteristic polynomial.
  \item Count the classes of each type and conclude:
        $GL_2(\mathbb F_q)$ has exactly $q^2 - 1$ conjugacy
        classes. \emph{(Count monic irreducible quadratics over
        $\mathbb F_q$; unordered pairs $\{a, b\}$; remember
        invertibility constrains constant terms.)}
  \item (Cyclic is generic) Show that $M \in M_2(\mathbb F_q)$
        fails to be cyclic iff $M$ is scalar, and deduce that a
        uniformly random $2\times2$ matrix over $\mathbb F_q$ is
        cyclic with probability $1 - q^{-3}$. State the analogous
        heuristic for $M_n$ and large $q$ (no proof required):
        non-cyclic matrices are rare --- which is why
        \cref{pb:b3:modules:1}'s Part III needed real work only
        past the generic case.
\end{enumerate}

\medskip
\noindent\textbf{Part V --- Complements.}
\begin{enumerate}[resume]
  \item (Center of the commutant) Show that the center of the
        algebra $\mathcal C(u)$ is exactly $K[u]$ \emph{(combine
        the two inclusions of Part III)}. Deduce that
        $\mathcal C(u)$ is commutative iff $u$ is cyclic ---
        recovering the equality case of question 8 by a purely
        structural route.
  \item (Which dimensions occur?) Deduce from Frobenius'
        formula that $\dim\mathcal C(u) \equiv n \pmod 2$ for
        every $u$. Then determine the exact set of values taken
        by $\dim \mathcal C(u)$ as $u$ ranges over
        $\mathcal L(V)$ with $\dim V = 4$: show it is $\{4, 6,
        8, 10, 16\}$ \emph{(enumerate the degree sequences $n_1
        \leq \dots \leq n_s$ summing to $4$ and realize each by
        a nilpotent)}. In particular $12$ and $14$, though of
        the right parity, are not attained: the parity
        constraint is necessary but not sufficient.
  \item (Class equation of $GL_2(\mathbb F_3)$) For $q = 3$,
        compute the size of each conjugacy class of question 20
        via orbit--stabilizer: the centralizer of a
        \emph{cyclic} $M$ in $GL_2(\mathbb F_q)$ is the unit
        group of $K[M]$ (question 11). Identify $K[M]$ in the
        three noncentral types, list the three monic irreducible
        quadratics over $\mathbb F_3$, and verify the class
        equation
        \[
        48 = \abs{GL_2(\mathbb F_3)}
        = 2\cdot1 + 1\cdot12 + 2\cdot8 + 3\cdot6,
        \]
        with $2 + 1 + 2 + 3 = 8 = q^2 - 1$ classes, as
        predicted by question 21.
\end{enumerate}
\end{problem}
